
\section*{Аннотация}

Работа посвящена разработке новых эвристических методов планирования запросов в реляционных СУБД. В литературе известно достаточно много алгоритмов планирования, 
но в индустриальных применениях известны лишь базовые подходы, опирающиеся на динамическое программирование, жадные и генетические алгоритмы и их простейшие 
эвристические комбинации, позволяющие выбирать тот или иной алгоритм в зависимости от числа таблиц в запросе. Данная работа преследует цель получить глубокое понимание 
того, как разные алгоритмы возможно комбинировать в ходе планирования запроса, делая переключение между алгоритмами в зависимости от топологии и сложности графа запроса. Ожидаемым результатом работы является решение оптимизационной задачи планирования с помощью новых эвристик, реализация решения в 
виде модификации планировщика в ядре реляционной СУБД с открытым кодом и экспериментальная оценка на известных индустриальных бенчмарках.

\subsection*{Постановка задачи}
\textbf{Дано}:
\begin{itemize}
    \item Реляционная СУБД PostgreSQL вместе с её планировщиком и стоимостной моделью;
    \item SQL-запросы индустриальных бенчмарков;
    \item Статистики таблиц( число строк, гистограммы значений, количество различных значений, наличие индексов);
    \item Аппаратные ресурсы исполнения, отражённые в параметрах стоимостной модели.
    \item Данные, на которых выполняются запросы. Представлены таблицами.
\end{itemize}

\textbf{Требуется получить}:
\begin{itemize}
    \item Новый метод выбора порядка соединений, реализованный в ядре
    PostgreSQL и интегрированный в конвейер планирования;
    \item Экспериментальная оценка на индустриальных бенчмарках;
    \item Сравнение нового подхода со стандартным динамическим программированием PostgreSQL,
    а также с известными альтернативами.
\end{itemize}